\documentclass[12pt,a4paper]{article}

% Paquetes esenciales
\usepackage[spanish]{babel}
\usepackage[utf8]{inputenc}
\usepackage[T1]{fontenc}
\usepackage{geometry}
\geometry{a4paper, margin=2.5cm}
\usepackage{hyperref}
\usepackage{titlesec}
\usepackage{enumitem}
\usepackage{graphicx}
\usepackage{xcolor}

% Configuración de colores y enlaces
\definecolor{darkgreen}{RGB}{34, 139, 34}
\hypersetup{
    colorlinks=true,
    linkcolor=darkgreen,
    urlcolor=darkgreen
}

% Formato de títulos
\titleformat{\section}{\Large\bfseries\color{darkgreen}}{}{0em}{}[\titlerule]
\titleformat{\subsection}{\large\bfseries}{}{0em}{}

\begin{document}

\begin{center}
    {\Huge \textbf{\color{darkgreen} Agente IA Agro-Inteligente}} \\[0.5cm]
    {\Large \textbf{Resumen del Proyecto y Requerimientos Tecnológicos}} \\[1cm]
    {\large \textbf{Autor:} Prof. César Rodríguez, C.I 5701716} \\
    {\small Dpto. de Física, Universidad de Oriente} \\
    {\small Núcleo de Sucre, Venezuela} \\[0.5cm]
    {\large Preparado para: Evaluación de Inversión y Financiación} \\[0.5cm]
    {\large \today}
\end{center}

\vspace{1cm}

\section{Resumen Ejecutivo}
El proyecto \textbf{Agente IA Agro-Inteligente} propone una solución vanguardista para la automatización y optimización de procesos agrícolas. Su núcleo radica en una arquitectura de software de inteligencia artificial de tres capas que separa la toma de decisiones probabilística de la ejecución determinista. Esto permite que el sistema analice datos en tiempo real, tome decisiones informadas y ejecute acciones sobre el terreno con un nivel de fiabilidad superior al 90\%, eliminando los errores comunes en los sistemas de IA tradicionales.

A través de su capacidad de ``autocuración'' (Self-healing) y memoria evolutiva, el agente no solo opera la maquinaria y los sensores agrícolas, sino que aprende de las anomalías del entorno, adaptando sus manuales de operación de forma autónoma. Además, el ecosistema integra tecnología Blockchain bajo el modelo DePIN, incentivando el despliegue de hardware físico y la recolección de datos valiosos mediante un token nativo, abriendo nuevas vías de monetización y financiamiento.

\section{Arquitectura Tecnológica (Ventaja Competitiva)}
El sistema se desmarca de las soluciones de mercado convencionales al implementar un modelo altamente escalable:
\begin{itemize}[label=$\bullet$]
    \item \textbf{Capa 1 (Directivas):} Manuales de operación estandarizados que definen los procesos agrícolas óptimos.
    \item \textbf{Capa 2 (Orquestación Inteligente):} Un motor cognitivo que enruta tareas, valida entradas de sensores y gestiona la recuperación ante errores.
    \item \textbf{Capa 3 (Ejecución):} Scripts de control deterministas aislados en contenedores seguros, garantizando que el hardware reaccione sin demoras ni fallos de lógica.
\end{itemize}

\section{Requerimientos de Hardware}
Para llevar a cabo el despliegue del proyecto y garantizar su máxima eficiencia, se requiere un conjunto de hardware especializado, dividido en el procesamiento central (Edge/Local) y la red de sensores en campo:

\subsection{1. Nodos de Campo y Telemetría (Microcontroladores)}
La recolección de datos y el control directo de actuadores (como sistemas de riego o ventilación) se basa en hardware de grado industrial pero de bajo costo:
\begin{itemize}
    \item \textbf{Módulos ESP32-S3:} El núcleo del hardware de campo. Seleccionado por su capacidad de procesamiento para IA en el borde (Edge AI), soporte nativo de Wi-Fi/Bluetooth para la topología de red, y su alta compatibilidad con nuestros entornos en \textit{C/C++ Sandbox}.
    \item \textbf{Sensores Agrícolas de Precisión:} 
    \begin{itemize}
        \item Sensores de humedad y temperatura del suelo.
        \item Estaciones meteorológicas compactas (anemómetros, pluviómetros).
        \item Sensores de variables ambientales (CO2, radiación solar).
    \end{itemize}
    \item \textbf{Actuadores:} Electroválvulas y relés de estado sólido operados directamente por los ESP32-S3 para automatizar el fertirriego y otras labores.
\end{itemize}

\subsection{2. Servidor de Orquestación (Edge/Local Processing)}
El 'cerebro' del sistema opera a nivel local para garantizar que la granja funcione incluso sin conexión a internet, requiriendo un servidor o pasarela (Gateway) con las siguientes características:
\begin{itemize}
    \item \textbf{Procesamiento:} CPU basada en arquitectura x86\_64 (compatible con sistemas base Linux Mint/Ubuntu).
    \item \textbf{Memoria RAM:} Mínimo \textbf{32GB de RAM} dedicados estrictamente a la capa de orquestación, entornos aislados (Conda), bases de datos vectoriales (ChromaDB) y la carga en memoria de los modelos de inteligencia artificial.
    \item \textbf{Procesamiento Gráfico (GPU):} Tarjeta gráfica dedicada para acelerar el procesamiento de datos agrícolas en tiempo real y permitir la ejecución de Modelos de Lenguaje Grandes (LLMs) de forma local, garantizando mayor privacidad y autonomía total sin dependencia de la nube.
    \item \textbf{Almacenamiento:} Disco de estado sólido (SSD) para el alojamiento del historial de métricas, contenedores Docker y el registro del estado de ejecución.
    \item \textbf{Red y Conectividad:} Routers industriales para gestionar la red local de los microcontroladores ESP32 y conexión hacia la nube para actualizaciones de los modelos de IA.
\end{itemize}

\section{Viabilidad y Retorno de Inversión (ROI)}
Al invertir en el \textbf{Agente IA Agro-Inteligente}, se está apostando por una plataforma que reduce drásticamente los costos operativos. La optimización en el uso del agua y fertilizantes, combinada con la reducción de la intervención humana mediante un sistema informático que \textit{se diagnostica y repara a sí mismo}, asegura un rápido retorno de inversión. 

La separación de la IA probabilística en la nube del hardware determinista en el campo (ESP32-S3) protege la producción contra fallas sistémicas, asegurando una agricultura de precisión escalable y rentable.

\section{Integración Web3 y Economía del Token (DePIN)}
Para financiar la expansión del proyecto e incentivar la adopción masiva, el sistema integra una capa económica basada en la tecnología Blockchain, alineándose con las tendencias de \textbf{DePIN} (Redes de Infraestructura Física Descentralizada) y \textbf{ReFi} (Finanzas Regenerativas).

Se introduce el \textbf{Agro IA Token (AGRO)}, un criptoactivo (Smart Contract ERC-20) diseñado con utilidades claras dentro del ecosistema:
\begin{itemize}[label=$\bullet$]
    \item \textbf{Incentivo por Datos (Minteo):} Los agricultores que desplieguen Nodos ESP32-S3 y compartan datos de telemetría (anonimizados) con la red global de investigación reciben tokens AGRO como recompensa.
    \item \textbf{Acceso a Orquestación Avanzada:} El pago por actualizaciones de los modelos de IA predictiva (Capa 2) y licencias de uso comercial del software se realiza mediante el token.
    \item \textbf{Mecanismo Deflacionario (Quema):} Una porción de los tokens utilizados para pagar licencias es "quemada" (destruida permanentemente a través del contrato inteligente), reduciendo la oferta total de tokens y protegiendo su valor en el mercado para los inversores tempranos.
\end{itemize}

Esta estructura dual convierte a cada granja no solo en un consumidor de tecnología robótica, sino en un nodo generador de ingresos pasivos dentro de una red global descentralizada de inteligencia agrícola.

\section{Glosario de Términos Técnicos}
\begin{itemize}[label=$\bullet$]
    \item \textbf{ESP32-S3:} Microcontrolador industrial de bajo costo y alta eficiencia con conectividad Wi-Fi/Bluetooth, diseñado específicamente para ejecutar algoritmos de Inteligencia Artificial directamente en dispositivos físicos.
    \item \textbf{GPU (Graphics Processing Unit):} Unidad de Procesamiento Gráfico. Hardware especializado en realizar operaciones matemáticas complejas en paralelo, indispensable para acelerar el funcionamiento de la Inteligencia Artificial de forma local.
    \item \textbf{LLM (Large Language Model):} Modelo de Inteligencia Artificial avanzado que comprende y razona utilizando lenguaje natural, actuando como el motor cognitivo de toma de decisiones.
    \item \textbf{Edge AI (IA en el Borde):} Procesamiento de datos con inteligencia artificial directamente en el lugar donde se generan (el servidor local de la granja), sin depender de la nube. Asegura respuestas inmediatas y funcionamiento sin internet.
    \item \textbf{ChromaDB / Base de Datos Vectorial:} Sistema de almacenamiento estructurado para IA que funciona como una "memoria a largo plazo", permitiendo al agente recordar configuraciones, errores pasados y aprendizajes.
    \item \textbf{Docker / Contenedores:} Tecnología que empaqueta el software y sus dependencias en entornos cerrados, garantizando que el código de ejecución no sufra fallos por incompatibilidades del sistema operativo.
    \item \textbf{Determinista vs. Probabilística:} La ejecución \textit{determinista} asegura que una orden siempre produzca el mismo resultado mecánico exacto. La IA es \textit{probabilística} porque deduce y puede variar su respuesta. El sistema es seguro porque combina ambas aislando el razonamiento de la acción física mecánica.
    \item \textbf{DePIN:} \textit{Decentralized Physical Infrastructure Networks}. Proyectos que utilizan tokens en la blockchain para incentivar a comunidades a desplegar y mantener hardware físico, como nuestros nodos agrícolas.
    \item \textbf{Token AGRO (ERC-20):} Criptomoneda nativa del proyecto desplegada mediante un Contrato Inteligente, utilizada para la gobernanza, recompensas e intercambio de valor en el ecosistema.
\end{itemize}

\end{document}